\section{Introduction}\label{intro}

Supervised Fine-Tuning (SFT) is a widely used post-training method in which the pretrained model
is adapted using labelled input-output examples. In visual question answering (VQA) tasks, this 
typically involves conditioning a vision-language model (VLM) on an image and question and 
optimising it to reproduce the ground-truth answer. Albeit effective, SFT learns from fixed
reference outputs during training, whereas at inference time the model must generate and condition 
on its own predictions. This difference between training trajectories and model-generated inference
trajectories can introduce a distribution mismatch.

On-policy distillation (OPD) addresses this issue by allowing a student model to generate its
own trajectories while a teacher model provides dense token-level supervision on those same 
trajectories. However, conventional OPD normally relies on a separate, stronger teacher model which
can increase computation and memory requirements. To address both the SFT and the conventional OPD
limitations, \textcite{zhao2601self} introduced On-Policy Self-Distillation (OPSD). OPSD is a fine-tuning
technique in which a single model acts as both student and teacher under different contexts. The student
would generate using standard inference context, while the teacher would receive privileged information
that is only available during training.

Majority of OPSD research has primarily focused on language-model reasoning. More recently,
multimodal work such as Vision-OPD has demonstrated that on-policy distillation can also be applied
to VLMs by providing a cropped, evidence-specific image to the teacher while the student receives the
full image \parencite{yuan2026vision}. This provides a natural form of privileged visual information since the teacher receives
a more focused representation of the evidence relevant to the question.

This work investigates this principle for VQA using the GQA dataset. GQA provides structured 
scene graphs and question grounding information that enable deterministic construction of
question-relevant image crops without using the ground-truth answer \parencite{hudson2019gqa}. It is based on the Visual Genome 
dataset \parencite{krishna2017visual}. 

The approach is evaluated across three VLMs with different architectures and capacities: SmolVLM-500M-Instruct,
BLIP, and MiniCPM-V 4.6. Each model is assessed under three conditions: the pretrained baseline,
conventional SFT and OPSD. This enables a controlled experiment in which direct supervised adaptation can
be compared to privileged-context self-distillation.

The study addresses the following research questions:

\begin{itemize}
\item \textbf{RQ1:} To what extent does OPSD improve VQA performance compared with conventional SFT?
\item \textbf{RQ2:} How does the relative OPSD-over-SFT performance vary across VLMs with different architectures and model characteristics?
\end{itemize}

The main contribution of this work is a cross-model empirical evaluation of OPSD for VQA using
question-grounded visual crops as privileged information.